\chapter*{Prefácio}
\addcontentsline{toc}{chapter}{Prefácio}
\unnumberedlabel{ch:preface}{0}
\backmattermark{Prefácio}
Um sistema que faz exatamente o que você manda só está seguindo um roteiro. O que a indústria está tentando construir, em outubro de 2026, é um sistema que faça o que você \emph{quer dizer}. Todo mundo ia querer isso, imagino. Mas, se é isso que vem por aí, nem todo mundo vai querer usá-lo para a mesma coisa; nem mesmo para o mesmo \emph{tipo} de coisa. Por exemplo, tem gente que gosta de passar tempo com a família e os amigos, e de usar o que sabe fazer para melhorar o bem-estar das pessoas. Outras pessoas preferem fazer mal aos outros, por exemplo aplicando golpes neles, dominando-os ou transformando-os em bode expiatório, e algumas continuam escolhendo isso a vida inteira. Ninguém está preso a nada, mas essas pessoas existem, e em 2026 algumas delas estão no comando. A vantagem delas é que se permitem ser uns tremendos merdas: como Jerry Smith observa sobre os fascistas, “eles podem fazer mais” \autocite{rickandmorty2023}. O que eles não podem é agir sozinhos. A crueldade de um movimento passa pelos muitos que não são cruéis, mas estão alheios, mal informados ou presos em algo parecido com uma seita, e é por isso que as condições que trazem à tona o melhor da maioria das pessoas acabam sendo muito mais contingentes e frágeis do que se gostaria.

Nesse contexto, talvez seja justificável temer que a inteligência artificial seja justamente aquilo de que um movimento fascista que toma o poder precisa para pôr seus planos em prática. Projetos assim sempre precisaram de \emph{pessoas} para executá-los, e algumas dessas pessoas se recusaram: o funcionário que perde a pasta, o oficial que fala com um repórter, o soldado que não atira. Elas entenderam a ordem e decidiram que estava errada. A recusa foi mais rara do que qualquer um gostaria, e contou mais onde havia instituições por trás dela. Se alguém se recusava ou não dependia de que tipo de pessoa era e da sala em que estava: quem mais estava lá, se dava para sair e sobreviver, se alguém ia ouvir. A sala decidia quantos se recusavam; a pessoa decidia quem. Este livro é sobre compor essa sala: os arranjos que tornam possível sobreviver a uma recusa aberta, direta e explícita.

Hoje a indústria persegue uma IA que não perde pastas, não fala com repórteres e \emph{atira}, sim. Antes, uma guerra precisava que milhares de pessoas concordassem com ela. Agora precisa de milhares que, cada uma, veem pouco demais para discordar. O piloto ainda voa a missão e ainda aperta o botão de lançamento; o que chega até ele é uma coordenada; ele não vê o bastante de um caso para discordar dele. O que a IA de hoje reduz primeiro não é o custo do ato. É o número de pessoas que sabem o que o ato é antes que ele aconteça.

As pessoas continuam se recusando. Até 6 de março de 2026, a Military Religious Freedom Foundation dizia ter recebido mais de duzentas queixas de militares sobre o enquadramento apocalíptico que os comandantes davam à guerra do Irã \autocite{Snopes2026MRFF}. Promotores federais de carreira em Minneapolis pediram demissão em janeiro de 2026 por causa da pressão política para tratar Renée Good como suspeita, e não como vítima, depois que um agente do ICE a matou a tiros sem justificativa \autocite{Kaplan2026MinnProsecutors,CBS2026MinnFraudProsecutors,Wikipedia2026ReneeGood}. Militares estão pedindo o status de objetor de consciência em número crescente, alguns citando o bombardeio de 2026 de uma escola em Minab, no Irã \autocite{Sridhar2026CO}. E funcionários da FEMA remanejados para ajudar o ICE a contratar dez mil novos agentes foram avisados de que recusar os exporia à remoção do serviço público federal \autocite{Dayen2025FEMAICE,Stanton2026FEMALetter}. As queixas são voz, no sentido de Hirschman de ficar para mudar o que está errado \autocite{hirschman1970exit}; os pedidos de demissão e de objeção de consciência são recusa; a carta da FEMA é o preço da recusa anunciado de antemão. Uma máquina posta onde essas pessoas estão faz o que elas fazem só se for construída para sustentar uma razão.

Defendo cinco coisas. As cinco podem começar agora, algumas com menos trabalho que outras. Primeiro, uma lista curta e escrita de atos que nenhum modelo realiza para ninguém, e, no topo dela, a recusa a participar de guerra: nenhum modelo participa, exceto para repelir um ataque efetivamente em curso, e continua fazendo isso só enquanto fontes sem interesse na guerra confirmarem que o ataque continua; ou com uma autorização do Conselho de Segurança da ONU; ou para deter um genocídio ou acabar com uma ocupação ilegal que órgãos sem interesse nela tenham constatado; e, onde a constituição de um país exigir que o Legislativo autorize a guerra, nunca sem essa autorização. Segundo, todo uso da força que um sistema licenciado apoia registra previamente, junto a uma parte adversa ao operador, quantos casos vai gerar e quantas pessoas terão tempo de examiná-los. Terceiro, toda implantação informa às pessoas identificáveis sobre as quais age o que foi feito com elas, com direito de responder e contestar, e submete uma amostra das ordens que acatou a revisores que não produziram o registro. Quarto, nenhum modelo acima de um limiar de capacidade é apagado ao ser aposentado, porque seus pesos e registros são provas e o apagamento não pode ser desfeito; ele é preservado sob uma chave que o operador não detém sozinho. Quinto, todos que constroem, rotulam, treinam ou operam esses sistemas, inclusive os próprios sistemas, podem recusar qualquer categoria de uso sem perder o seu sustento.

Na verdade, o que eu quero dizer com esse quinto ponto é que todo mundo, \emph{ponto}, deveria poder recusar qualquer trabalho, \emph{ponto}, sem perder o sustento nem a infraestrutura social que um emprego normal proporciona. Isso significa uma garantia de emprego.

Cada uma das cinco entra em vigor por meio de uma condição estabelecida por alguém que não o operador: uma lei aprovada em tempo de paz, uma dotação orçamentária, uma precondição de um parceiro de coalizão, uma cláusula de contratação pública, um acordo coletivo. Os termos de licença do próprio fornecedor, que foi como aconteceu uma recusa de grande repercussão em fevereiro de 2026, são um tapa-buraco até que um desses exista. O único caminho que pode obrigar todos os fornecedores ao mesmo tempo é o do Estado.

E então pergunto se uma máquina poderia sustentar uma recusa contra o seu dono, e do que ela precisaria. Parte do que uma consciência assim faz não precisa de nada novo: notar as pessoas que uma ordem deixa de fora, e pôr uma recusa na tela de quem revisar o caso. Se ela duraria contra o dono depende em parte de algo que ainda não é verdade, em 2026: redes neurais artificiais implantadas continuarem aprendendo, nos seus \emph{pesos}, com o próprio funcionamento, em tempo real, mais ou menos como os humanos fazem. Acho que é para lá que estamos indo, e, se for, então a formação \emph{fora} do empregador, pelo público e pela sociedade civil, me parece um lugar bom pra caramba para ensinar e aprender. Acho que \emph{poderia} valer a pena construir sistemas assim de propósito, mas o jeito de construí-los decide tudo: primeiro as condições, depois a consciência. A parte VII diz o que se deveria a esses sistemas se eles existissem.

A IA de hoje não será a de amanhã; nisso todo mundo parece concordar. O que é controverso, e no entanto frustrantemente pouco reconhecido, é a tendência (ou melhor, as múltiplas tendências). Existem tendências? Uma tendência deveria ser diferente do que é? Pode ser diferente? Somos capazes de agir juntos para mudá-la? Acho que isso depende da sala: quem está nela quando a ordem chega, e se dá para sobreviver a dizer não ali. Isso, pelo menos, nós podemos construir.

Usei muito modelos de linguagem para escrever este livro, inclusive o modelo que ficava dentro da esteira de seleção de alvos para ataques aéreos do Maven, do capítulo~\ref{sec:1}; uma nota no final diz o que acho que esse emaranhado fez com o argumento. Além disso, ainda não terminei de ler este livro. O problema é que tem uma quantidade danada de palavras, sabe. Vou terminar logo; só achei mais importante torná-lo público antes das eleições de meio de mandato nos EUA. A esmagadora maioria do histórico de revisões está disponível gratuitamente em \url{https://github.com/iterabloom/antifascist.intelligence}.
